% 補足資料 S1 の表：ネットワークの各部の大きさと学習の設定（第7弾，2026-10-28）。
% 出典：config/model/SuperpositionNetworkMotionPleasureHold{,Other}/default.yml，.../OneBit/onebit.yml，model/modules.py・model.py・model_v12.py・model_v14.py，
%   exp/runner.py，config/exp の段階I（MSE・L1）と段階IIの設定，simulation/collect_data_v11.py，sim_v10/config/agent，analyze/v12/eval_g5_v12.py・eval_trial_v12.py，
%   analyze/v11/localise_v11.py・eval_trial_v11.py（boot_ci），analyze/v14/aggregate_inverse_v14.py。層の出力の大きさはモデルを CPU で組み立てて確かめた。
%   照合の記録：paper/check/numbers_check_v7.md §1。
\begin{tabular}{@{}p{0.15\linewidth}p{0.22\linewidth}p{0.56\linewidth}@{}}
\toprule
部分 & 設定 & 値 \\
\midrule
視覚エンコーダ & 入力 & 16×64×3（A-1 の視界，RGB） \\
（1・2 は同じ作り） & 畳み込み層1〜3 & 16・32・64 チャンネル，カーネル 3×10（縦×横），ストライド 2，パディング 1（出力 8×29，4×11，2×2），各層の後に ReLU \\
 & 全結合層 & 256 → 64，層正規化（2つのエンコーダで共有）の後に ReLU \\
 & 重み & 基準エンコーダに固定（すべての条件・期間で凍結） \\
\midrule
共有モジュール & LSTM & 入力 75（視覚 64 ＋ $x_t$ 11），隠れ状態 128．process-1・2 で重みを共有し，状態は別 \\
 & process-1 の入力 $x^1_t$ & 11 次元（運動2 ＋ 快・不快8 ＋ 隠れフラグ1） \\
 & process-2 の入力 $x^2_t$ & 段階I：11 次元すべて0．段階II：$[\,0,\,Q2_t,\,1\,]$（運動の2次元は0） \\
\midrule
視覚予測器 & 統合層 & $[h^1_t,h^2_t]$ の 256 → 64，ReLU．学習中は $h^1_t$，$h^2_t$ に確率 0.5 のドロップアウト \\
 & 復号部の全結合層 & 64 → 256，ReLU の後に 64×2×2 に並べ替え \\
 & 転置畳み込み層1〜3 & 32・16・3 チャンネル，カーネル 4×10・4×12・4×12，ストライド 2，パディング 1（出力 4×10，8×28，16×64），層1・2 の後に ReLU，層3 の後に tanh \\
\midrule
特徴予測器 & 入力 & $h^1_t$ または $h^2_t$（128．段階Iでは勾配を止める） \\
 & 全結合層 & 128 → 128 → 64，各層で層正規化の後に ReLU \\
 & 目標と損失 & 次の視界を視覚エンコーダ1（$h^1$ の側）・2（$h^2$ の側）に通した出力，二乗誤差 \\
\midrule
他者用モジュール & 入力 & マスク前の視覚エンコーダ2の出力（64） \\
 & LSTMCell，全結合層 & 隠れ状態 128，tanh を通した出力は 8 次元（1ビットは 1 次元） \\
 & 1ビットの形 $c^{+}$ & $Q2_t=u_t\,c^{+}$，$c^{+}=(-0.104$，$-0.127$，0，0，0，0，$+0.127$，$+0.104)$ \\
\midrule
段階Iの学習 & MSE 期間 & 200 エポック（毎エポック保存）．次の視界と特徴予測の二乗誤差で，共有モジュール・統合層・復号部・特徴予測器を学習 \\
 & L1 期間 & 200 エポック（10 エポックごとに保存）．MSE 期間の 200 から，次の視界の絶対誤差で復号部だけを学習 \\
 & 視覚のマスク & $t\ge1$ の各ステップで確率 0.99（$t=0$ は見せる） \\
 & 運動を隠す区間 & 長さ 5〜15 ステップの見える区間（$t=0$ から）と隠れた区間を交互に置く \\
\midrule
段階IIの学習 & 学習する部分 & 他者用モジュールだけ．段階Iの L1 期間の 200 から 200 エポック（10 エポックごとに保存） \\
 & 損失 & A-1 の次の視界の絶対誤差 ＋ 特徴予測の二乗誤差 \\
 & 視覚と運動 & $t\ge1$ で process-1・2 とも視覚を隠す（確率 1）．process-1 の運動は隠さない \\
\midrule
最適化（段階I・II） & AdamW，バッチ & 学習率 $10^{-3}$，重み減衰 0.01（バイアスは0），バッチ 10 エピソード \\
\midrule
データ & エピソード & 101 フレーム．各条件・各データで学習用 3000，評価用 300 \\
 & 休止 & 動く一歩ごとに確率 0.005 で始まり，10〜30 ステップ続く \\
\midrule
評価 & 自分側の積み上がり $g$ & $t=20$ から 15 ステップ止め，$d=0,\dots,10$ 歩動いた視界と比べる \\
 & G5b & $a=0.28$，$t=20$〜29 の 10 ステップ \\
 & 照合 & 目標のフレーム数 $T=50$（主） \\
 & 信頼区間 & エピソード単位のブートストラップ 1000 回，95\%（パーセンタイル） \\
\bottomrule
\end{tabular}
